\chapter{Process Flow: Cue Typing, Spreading Activation, and Pattern Completion}

\section{The User Types an Instruction, Not a Question}
\textbf{Trigger:} The user types \emph{``implement feature A''} in the Tottie chat (or a coding agent
issues the same over MCP).

\textbf{Action / UI State:} The sprite enters \texttt{thinking}; the Memory Panel opens a
``RECALL'' block with a \emph{cue} chip showing the classified type.

\textbf{System Logic:} The \textbf{Cue Classifier} labels the turn \texttt{COMMAND/INTENT} and extracts
entities and the project scope (\texttt{\#project/A}). The retrieval \textbf{policy} for a command is
not ``find similar text''; it is \emph{assemble what a competent person needs to act}: project
bindings, conventions, prior decisions on \texttt{A}, lessons tied to \texttt{A}, and constraints.

\textbf{Data Flow \& Mutators:} A retrieval request is issued with the intent, entities, and scope; a
\texttt{memory\_event} of kind \texttt{recall} begins.

\section{The Deterministic Rules (implementation contract)}
Rules run in order; STATEMENT is the default, so every turn classifies with zero model calls.
An LLM fallback for genuinely ambiguous turns is deferred past the hackathon.
\begin{longtable}{@{}p{0.20\textwidth}p{0.68\textwidth}@{}}
\toprule
Cue type & Rule (first match wins) \\
\midrule
CORRECTION & contains \texttt{actually}, \texttt{no wait}, \texttt{correction}, \texttt{I meant} \\
REFERENCE & contains \texttt{that} + noun, \texttt{the same}, \texttt{we discussed}, \texttt{you mentioned} \\
QUESTION & ends with \texttt{?}, or starts with who/what/when/where/why/how/which \\
COMMAND & starts with an imperative verb (implement, add, fix, create, update, refactor, build, write, run, show, remind) \\
CHITCHAT & matches \texttt{\^{}(thanks|thank you|hi|hello|hey|bye|ok|okay|cool|great)[!.]*\$} (case-insensitive) \\
STATEMENT & default: anything carrying new information \\
\bottomrule
\end{longtable}
A turn matching both CORRECTION and another type is a CORRECTION (the reconsolidation route
depends on it). ``Remind me to\ldots'' is a COMMAND that additionally registers an intention
(see \emph{The Recall Path}: the prospective check). Entities for seeding and event\_cue
overlap come from rules only: \texttt{\#tags}, \texttt{@mentions}, and quoted strings. No NER
model in the hackathon build.

\section{Seeding the Associative Network}
\textbf{Trigger:} The classified cue.
\textbf{Action / UI State:} The Panel shows seed chips (e.g.\ \texttt{\#project/A}, \texttt{feature A}).
\textbf{System Logic:} Nodes are seeded directly from the parsed entities (bypassing similarity), and
the three legs run in parallel: \textbf{semantic} (embedding nearest neighbours), \textbf{lexical}
(full-text match), and \textbf{associative} (graph nodes for the seeded tags/entities).
\textbf{Data Flow \& Mutators:} A candidate set of memory ids with initial activation.

\section{Spreading Activation}
\textbf{Trigger:} The seeded candidate set.
\textbf{Action / UI State:} The Panel renders each hop as it happens: ``spread 1 hop (w = 0.6)''.
\textbf{System Logic:} The \textbf{Activator} walks \texttt{memory\_links} outward from the seeds,
propagating activation with a per-hop decay and a threshold cut-off. Linked memories that were never
textually similar to the instruction still light up --- because they are \emph{associated}, exactly as
a band name revives an unrelated scene from fifteen years ago.
\textbf{Data Flow \& Mutators:} A ranked activation vector over memory ids.

\section{Pattern Completion}
\textbf{Trigger:} A high-activation region of the graph.
\textbf{Action / UI State:} The Panel groups the returned memories into one coherent cluster.
\textbf{System Logic:} The \textbf{Completer} takes the top-activated \emph{cluster} and reconstructs a
coherent, self-consistent context --- the biological analogue of CA3 recurrent excitation restoring a
whole memory from a partial cue.
\textbf{Data Flow \& Mutators:} The cluster replaces the loose hit list.

\section{Curate and Inject}
\textbf{Trigger:} The completed cluster.
\textbf{Action / UI State:} The Panel's ``CONTEXT INJECTED THIS TURN'' block shows the final Context
Card, each line annotated with \emph{why it surfaced} (cue type $\rightarrow$ seed $\rightarrow$ hops).
\textbf{System Logic:} The \textbf{Curator} assembles the card within a token budget and appends the
``probe deeper'' pointers.
\textbf{Data Flow \& Mutators:} The card is injected into the prompt; a \texttt{recall} event is
streamed to the Panel.

\section{Temporal Scope: The ``What Did I Just Do'' Branch}
\textbf{Trigger:} A QUESTION carrying recency markers: \emph{just, just now, so far, last N,
recent, lately, where were we, recap, what have we done}.
\textbf{Action / UI State:} The Panel renders an ordered session timeline --- one line per turn,
oldest first --- instead of a ranked hit list.
\textbf{System Logic:} Similarity ranking is the wrong tool for a recency question (it returns
the most \emph{similar} turns, not the \emph{latest} ones). The Retriever therefore takes the
replay branch: it reads the session's \texttt{memory\_event} rows in monotonic-id order with
their linked episodic rows, and presents them as ordered chunks (one line per turn from the
stored \texttt{scribe\_action} summaries). This is machine \textbf{serial recall}: every served
memory gets its \texttt{access\_count} increment, so the recap doubles as \textbf{rehearsal on
access} (testing effect). Per the governing principle we copy the solution (ordered replay +
rehearsal) and refuse the constraints (no 15--30~second buffer, no transposition errors, no
output interference) --- the machine replay is exact.
\textbf{Data Flow \& Mutators:} An ordered read of \texttt{(session\_id, id)} plus
\texttt{access\_count} bumps; no new rows, no similarity search.

\section{Inhibition and Abstention}
\textbf{Trigger:} Nothing clears the activation floor, or everything retrieved is superseded.
\textbf{System Logic:} Inhibit the results and inject \emph{nothing}. Abstaining is a correct outcome;
injecting noise is not.
